\section{Adiabatic turbulent boundary layer profiles}
\label{sec:numerical-vv:adiabatic-turbulent-boundary-layer-profiles}

The Mach-4.5 flat-plate case investigated in this section assesses the predictions of the
\(k\)--\(\omega\) shear-stress-transport (SST) and Spalart--Allmaras (SA)
models in an attached, compressible turbulent boundary layer. The reference
solutions are the Wind-US calculations distributed through the NASA Glenn
NPARC Alliance Validation Archive~\citep{DippoldLenahan2012Mach45FlatPlate}.
The experimental basis is the flat-plate investigation of
Coles~\citep{Coles1953SupersonicFlatPlate}, with the experimental data used here
are taken from the NASA archive\footnote{\url{https://www.grc.nasa.gov/WWW/wind/valid/fpmach45/fpmach45.html} [accessed June 11th 2026]}. The principal objective is to compare
OpenFOAM/HiSA with NASA's implementation of the corresponding turbulence
models. This comparison checks whether the predicted velocity profiles and
wall shear reproduce the reference behaviour qualitatively. The measurements provide an
additional assessment against physical observations.
\par\medskip

The selected NASA solutions use Wind-US~3.162 on a structured mesh with its
unstructured solver, denoted \emph{Str-Uns} in the archive and figure legends.
Both reference calculations assume a fully turbulent boundary
layer. The comparison therefore concerns the development of an attached
turbulent layer, while the different transition and tripping histories in
the experiment remain relevant when interpreting agreement with Coles' experimental results.
The present section considers the OpenFOAM native implementation of \(k\)--\(\omega\) SST and the density-corrected SA implementation (see Section \ref{subsec:ext:spalart-allmaras-of-correction}).

\paragraph{General case setup.}

The flow is parallel to a flat-plate at \(M_\infty=4.512\), with no imposed
geometrical turning or external streamwise pressure gradient. The lower
boundary upstream of \(x=0\) is a slip wall, where the subsequent no-slip, adiabatic plate
begins at \(x=0\). No transition model or explicit trip is included, the no-slip wall is solved as fully turbulent. Both OpenFOAM/HiSA cases use one spanwise cell with \texttt{empty} side
boundaries. The inlet state is fixed to the prescribed freestream,
corresponding to the frozen inlet of the NASA case. The upper and rear boundaries use characteristic farfield conditions
(Section~\ref{subsec:numerical-framework:hisa}). At the viscous wall, velocity is zero and
pressure and temperature have zero normal gradient, so the adiabatic wall
temperature is obtained as part of the solution. The upstream slip surface
is excluded from the viscous-wall distance calculation.
\par\medskip

Air is represented by a calorically perfect ideal gas with
\(c_p=1004.5002\,\mathrm{J\,kg^{-1}\,K^{-1}}\),
\(R=287.00005\,\mathrm{J\,kg^{-1}\,K^{-1}}\), and \(\gamma=1.4\).
The molecular viscosity matches the Wind-US Sutherland law after conversion
to SI units,
\begin{equation}
    \mu(T)=A_s\frac{T^{3/2}}{T+T_s},
    \qquad A_s=1.49610554\times10^{-6}\,\mathrm{Pa\,s\,K^{-1/2}},
    \qquad T_s=120\,\mathrm{K}.
    \label{eq:coles-sutherland-law}
\end{equation}
A constant molecular Prandtl number is imposed through
\(\lambda=c_p\mu/Pr\), with \(Pr=0.72\). Both turbulence models use a constant \(Pr_t=0.9\), as documented in the NASA Wind-US user guide, thus a reasonable assumption to also have been used for the NASA reported results. Table~\ref{tab:coles-flow-conditions} summarizes the matched physical
conditions. The freestream viscosity is \(\mu_\infty=3.89802\times10^{-6}\,\mathrm{Pa\,s}\), giving
\(Re_x/x=6.98354\times10^6\,\mathrm{m^{-1}}\).

% Sources: WSL /home/mgnadlinger/Coles_TurbZPGFP/{kwSST,SA}/0/include/setup,
% constant/{thermophysicalProperties,turbulenceProperties}, and 20000 fields.
% NASA Downloads/Cases/Str-Uns_{SST,SA}/FP4p5.{dat,lis}.

\begin{table}[H]
    \centering
    \caption{Freestream and thermal-model quantities prescribed for the
        OpenFOAM/HiSA Coles flat-plate cases.}
    \label{tab:coles-flow-conditions}
    \small
    \setlength{\tabcolsep}{4pt}
    \renewcommand{\arraystretch}{1.12}
    \begin{tabular}{@{}p{0.62\linewidth}>{\centering\arraybackslash}p{0.31\linewidth}@{}}
        \toprule
        Quantity & {OpenFOAM/HiSA} \\
        \midrule
        Mach number \(M_\infty\) & \num{4.512} \\
        Static pressure \(p_\infty\) [\(\mathrm{Pa}\)]
            & \num{671.719} \\
        Static temperature \(T_\infty\) [\(\mathrm{K}\)]
            & \num{60.4672} \\
        Freestream velocity \(U_\infty\) [\(\mathrm{m\,s^{-1}}\)]
            & \num{703.289} \\
        Heat-capacity ratio \(\gamma\) & \num{1.4} \\
        \raggedright Unit Reynolds number \(Re_x/x\) [\(\mathrm{m^{-1}}\)]
            & \num{6.984e6} \\
        Molecular Prandtl number \(Pr\)
            & \num{0.72} \\
        Turbulent Prandtl number \(Pr_t\) & \num{0.9} \\
        Wall thermal condition & {Adiabatic} \\
        \bottomrule
    \end{tabular}
\end{table}

\paragraph{Mesh and turbulence-model settings.}

Both cases use the structured mesh taken from the NASA Str-Uns
flat-plate validation archive~\citep{DippoldLenahan2012Mach45FlatPlate}
and converted to OpenFOAM, with coordinates converted from feet to metres.
Figure~\ref{fig:coles-nasa-mesh} shows the mesh and boundary-condition
layout shown in the archive.

\begin{figure}[H]
    \centering
    % Figure provenance (audited 2026-10-19; supplied NASA illustration):
    %   Included figure: 06_mesh_NASA_archive_Figure_5.png in this Figures directory.
    %   Retained NASA source image:
    %     /home/mgnadlinger/Coles_TurbZPGFP/NASA Downloads/Grids/Str-Str/FP4p5_BCs.png
    %   Windows copy:
    %     /mnt/c/Users/MichaelGnadlinger/OneDrive - Hypersonica/Hypersonica Team - Documents/02 Technical/03 Flight physics/16 Theses/2026/01 Gnadl/Files/ValData/Tunel/03 TurbZPGFP/01 Coles/NASA Downloads/Grids/Str-Str/FP4p5_BCs.png
    %   The included and WSL source images are pixel-identical RGB 488 x 379.
    %     File hashes differ because the manuscript copy omits PNG metadata.
    %     Source SHA-256: 4cb241bda0aa47ca6a5063887f359de35687f1c7c584b889ed38bcdf6f2b0765.
    %   This is the NASA boundary-condition schematic stored under Str-Str,
    %     not the separate oblique Str-Uns grid overview FP4p5_StrUns.png.
    %   Actual simulation grid provenance: NASA Downloads/Grids/Str-Uns/
    %     FP4p5_StrUns.{cgd,gg}, with OpenFOAM conversion retained in
    %     openfoam_FP4p5_StrUns/{convert.py,conversion_metadata.json}; these
    %     files are present under the WSL Coles_TurbZPGFP root.
    %   NASA Tecplot source: Grids/Str-Str/FP4p5_BCs.lay in the same NASA
    %     Downloads root; it loads FP4p5_StrStr.cgd and contains the displayed
    %     boundary-condition annotations. Both files are retained locally.
    %     This supplied illustration has no thesis Python plotting script.
    % End figure provenance.
    \includegraphics[width=0.8\linewidth]{04_OpenFOAM_HiSA_VnV/Figures/06_mesh_NASA_archive_Figure_5.png}
    \caption{Structured mesh and boundary-condition layout for the
        Mach-4.5 flat-plate case. Taken from the NASA NPARC Alliance
        Validation Archive~\citep{DippoldLenahan2012Mach45FlatPlate}.}
    \label{fig:coles-nasa-mesh}
\end{figure}

The two models use the inlet quantities and near-wall conditions listed in
Table~\ref{tab:coles-turbulence-setup}. These values
match the NASA inlet settings.

% Inlet quantities: 0/include/setup; wall resolution: 20000/yPlus and
% postProcessing/yPlus/4000/yPlus.dat at iteration 20000.

\begin{table}[H]
    \centering
    \caption{Turbulence inputs and recorded wall resolution of the two
        OpenFOAM/HiSA cases. Inlet viscosity ratios refer to the molecular
        kinematic viscosity used in the inlet setup. Wall \(y^+\) statistics
        are from the saved iteration 20\,000.}
    \label{tab:coles-turbulence-setup}
    \small
    \setlength{\tabcolsep}{4pt}
    \renewcommand{\arraystretch}{1.12}
    \begin{tabular}{@{}p{0.42\linewidth}S[table-format=1.5e-1,table-column-width=0.25\linewidth] S[table-format=1.5e-1,table-column-width=0.26\linewidth]@{}}
        \toprule
        Quantity & {SST} & {SA} \\
        \midrule
        \(k_\infty\) [\(\mathrm{m^2\,s^{-2}}\)]
            & 2.32369e-3 & {---} \\
        \(\omega_\infty\) [\(\mathrm{s^{-1}}\)]
            & 2.30738e4 & {---} \\
        \(\widetilde{\nu}_\infty\) [\(\mathrm{m^2\,s^{-1}}\)]
            & {---} & 5.03533e-4 \\
        Inlet viscosity ratio
            & {\(\nu_{t,\infty}/\nu_\infty=0.001\)}
            & {\(\widetilde{\nu}_\infty/\nu_\infty=5\)} \\
        \bottomrule
    \end{tabular}
\end{table}

Iterative convergence is assessed from the stability of the physical wall shear
stress. Across the six saved solutions from iterations 18\,000 to
20\,000, the maximum peak-to-peak downstream skin-friction variation
for \(x\geq0.1\,\mathrm{m}\) is \(0.0875\%\) for SA and
\(0.00352\%\) for SST. At the profile station, the corresponding
variations are below \(0.0002\%\), while changes in the streamwise
velocity profiles are below \(0.00005\%\) of \(U_\infty\).
The remaining flow-residual jitter is treated as numerical; it does not
produce material changes in these comparison quantities. The figures
therefore use the accepted iteration-20\,000 solutions as fully converged.

\paragraph{Extraction and normalization.}

The streamwise and wall-normal velocities are denoted by \(u\) and \(v\) respectively.
Both are normalized with \(U_\infty=703.289\,\mathrm{m\,s^{-1}}\). The streamwise skin-friction
coefficient is evaluated as
\begin{equation}
    C_{f,x}(x)=\frac{|\tau_{w,x}(x)|}
        {\tfrac12\rho_\infty U_\infty^2},
    \qquad \rho_\infty=0.0387067\,\mathrm{kg\,m^{-3}}.
    \label{eq:coles-skin-friction}
\end{equation}
NASA lengths and velocities are converted from feet and feet per
second to SI units.

\subsection{Velocity profiles and skin-friction}
\label{subsec:coles-results}

Figure~\ref{fig:coles-velocity-profiles} compares both velocity components
at \(x=0.5461\,\mathrm{m}\). The streamwise profiles closely follow
the NASA solutions through the boundary layer and approach the same outer
velocity. the logarithmically plotted velocity profile in Figure~\ref{fig:coles-velocity-experiment} gives an even better view of the close matching between the numerical results by the NASA Wind-US Code and the herin used OpenFOAM/HiSA framework.

% Both Coles velocity plots are regenerated by Figures/plot_coles_profile_fonts.py.
% Inputs: Files/ValData/Tunel/03 TurbZPGFP/01 Coles/validation/
% outputs_nasa_matched_20000/data/{kwSST,SA}/profile_x0p5461.csv
% and the NASA Downloads reference tables; original simulations:
% /home/mgnadlinger/Coles_TurbZPGFP/{kwSST,SA}.
% Caption-sized typography update: curve data and axis limits preserved;
% source hashes are recorded in Figures/coles_profile_font_provenance.json.
\begin{figure}[H]
    \centering
    % Figure provenance (audited 2026-10-19; WSL simulations, Windows plotting inputs):
    %   Included figure: 06_Coles_velocity_profiles.pdf in this Figures directory.
    %   Current typography script:
    %     /mnt/c/Users/MichaelGnadlinger/OneDrive - Hypersonica/Hypersonica Team - Documents/02 Technical/03 Flight physics/16 Theses/2026/01 Gnadl/Docs_and_Presentations/Thesis_tex/thesis/04_OpenFOAM_HiSA_VnV/Figures/plot_coles_profile_fonts.py
    %   Original simulation root (WSL): /home/mgnadlinger/Coles_TurbZPGFP
    %     Cases: {kwSST,SA}/20000/U and constant/polyMesh; the extraction in
    %     validation/replot_nasa_matched.py interpolates cell-centre values in x
    %     to x = 0.5461 m and adds the no-slip wall point.
    %   Windows analysis / retained-data root:
    %     /mnt/c/Users/MichaelGnadlinger/OneDrive - Hypersonica/Hypersonica Team - Documents/02 Technical/03 Flight physics/16 Theses/2026/01 Gnadl/Files/ValData/Tunel/03 TurbZPGFP/01 Coles
    %   Inputs relative to that root:
    %     validation/outputs_nasa_matched_20000/data/{kwSST,SA}/profile_x0p5461.csv
    %     NASA Downloads/Cases/Str-Uns_{SST,SA}/FP4p5.BL.dat
    %     validation/outputs_nasa_matched_20000/validation_summary.json.
    %   Scientific plotting functions: validation/compare_nasa_figures.py;
    %     matching/extraction and preceding exports: validation/replot_nasa_matched.py.
    %   The current export changes typography/layout; curve coordinates and axis
    %     limits are checked against the original plotting functions by the script.
    %   Audit: coles_profile_font_provenance.json in this Figures directory.
    %     All 11 source hashes were independently rechecked and match. The raw
    %     WSL iteration-20000 U and wallShearStress hashes for both cases also
    %     match the accepted validation_summary.json. Normalization uses its
    %     U_inf = 703.2894492195271 m/s; NASA feet/ft/s are converted to SI.
    %   Plotting entry point: compare_nasa_figures.plot_fig6() (u and v).
    %   Preceding export: validation/outputs_nasa_matched_20000/fig6_velocity_profiles.pdf
    %     under either the Windows or WSL Coles root. The current manuscript file
    %     was subsequently re-exported by plot_coles_profile_fonts.py.
    % End figure provenance.
    \includegraphics[width=\linewidth]{04_OpenFOAM_HiSA_VnV/Figures/06_Coles_velocity_profiles.pdf}
    \caption[Coles case: velocity profiles compared with NASA Wind-US]{%
        Streamwise and wall-normal velocity profiles at
        \(x=21.5\,\mathrm{in}=0.5461\,\mathrm{m}\).
        Dark curves with symbols show present OpenFOAM/HiSA results; pale curves show the
        corresponding NASA Str-Uns SST and SA solutions.}
    \label{fig:coles-velocity-profiles}
\end{figure}

\begin{figure}[H]
    \centering
    % Figure provenance (audited 2026-10-19; WSL simulations, Windows plotting inputs):
    %   Included figure: 06_Coles_velocity_experiment.pdf in this Figures directory.
    %   Current typography script:
    %     /mnt/c/Users/MichaelGnadlinger/OneDrive - Hypersonica/Hypersonica Team - Documents/02 Technical/03 Flight physics/16 Theses/2026/01 Gnadl/Docs_and_Presentations/Thesis_tex/thesis/04_OpenFOAM_HiSA_VnV/Figures/plot_coles_profile_fonts.py
    %   Original simulation root (WSL): /home/mgnadlinger/Coles_TurbZPGFP
    %     Cases: {kwSST,SA}/20000/U and constant/polyMesh; the extraction in
    %     validation/replot_nasa_matched.py interpolates cell-centre values in x
    %     to x = 0.5461 m and adds the no-slip wall point.
    %   Windows analysis / retained-data root:
    %     /mnt/c/Users/MichaelGnadlinger/OneDrive - Hypersonica/Hypersonica Team - Documents/02 Technical/03 Flight physics/16 Theses/2026/01 Gnadl/Files/ValData/Tunel/03 TurbZPGFP/01 Coles
    %   Inputs relative to that root:
    %     validation/outputs_nasa_matched_20000/data/{kwSST,SA}/profile_x0p5461.csv
    %     NASA Downloads/Cases/Str-Uns_{SST,SA}/FP4p5.BL.dat
    %     validation/outputs_nasa_matched_20000/validation_summary.json.
    %   Scientific plotting functions: validation/compare_nasa_figures.py;
    %     matching/extraction and preceding exports: validation/replot_nasa_matched.py.
    %   The current export changes typography/layout; curve coordinates and axis
    %     limits are checked against the original plotting functions by the script.
    %   Audit: coles_profile_font_provenance.json in this Figures directory.
    %     All 11 source hashes were independently rechecked and match. The raw
    %     WSL iteration-20000 U and wallShearStress hashes for both cases also
    %     match the accepted validation_summary.json. Normalization uses its
    %     U_inf = 703.2894492195271 m/s; NASA feet/ft/s are converted to SI.
    %   Experimental input: NASA Downloads/Results/FP4p5-BL.dat, Experiment zone,
    %     relative to the Windows analysis root above; supplied Coles profile.
    %   Plotting entry point: compare_nasa_figures.plot_fig7() (u vs log wall distance).
    %   Preceding export: validation/outputs_nasa_matched_20000/fig7_u_vs_experiment.pdf
    %     under either the Windows or WSL Coles root. The current manuscript file
    %     was subsequently re-exported by plot_coles_profile_fonts.py.
    % End figure provenance.
    \includegraphics[width=0.85\linewidth]{04_OpenFOAM_HiSA_VnV/Figures/06_Coles_velocity_experiment.pdf}
    \caption[Coles case: streamwise velocity and experimental reference]{%
        Streamwise velocity at \(x=0.5461\,\mathrm{m}\), plotted against
        logarithmic wall distance. The experimental reference is the
        \emph{Experiment} profile supplied with the NASA archive; the
        numerical curves are the same SST and SA cases as in
        Figure~\ref{fig:coles-velocity-profiles}.}
    \label{fig:coles-velocity-experiment}
\end{figure}

Figure~\ref{fig:coles-skin-friction} presents the skin-friction coefficient,
including separate leading-edge and downstream views. Similarly to the velocity boundary layer profiles, the OpenFOAM/HiSA results match the NASA results, showing that shear-stress calculations are consistent between both codes. The three Coles fence-trip series provide experimental context for the downstream trend. 

\begin{figure}[H]
    \centering
    % Figure provenance (audited 2026-10-19; WSL original, retained Windows copy):
    %   Included figure: 06_Coles_skin_friction.pdf in this Figures directory.
    %   Exact original export (SHA-256 matched):
    %     /home/mgnadlinger/Coles_TurbZPGFP/validation/outputs_nasa_matched_20000/fig8_cfx.pdf
    %     fcbb0236d08c1ea329bcde05cf49880c32cce933a705ea6be865cce79b0498a1.
    %   Export / extraction script:
    %     /home/mgnadlinger/Coles_TurbZPGFP/validation/replot_nasa_matched.py
    %     (iteration 20000, accepted convergence; plot_fig8 from compare_nasa_figures.py).
    %   Windows copy of scripts, exported data and NASA reference tables:
    %     /mnt/c/Users/MichaelGnadlinger/OneDrive - Hypersonica/Hypersonica Team - Documents/02 Technical/03 Flight physics/16 Theses/2026/01 Gnadl/Files/ValData/Tunel/03 TurbZPGFP/01 Coles
    %   Data relative to that Windows root:
    %     validation/outputs_nasa_matched_20000/data/{kwSST,SA}/skin_friction.csv
    %     NASA Downloads/Cases/Str-Uns_{SST,SA}/FP4p5.Cf.dat
    %     NASA Downloads/Results/FP4p5-Cf.dat (three Coles fence-trip zones).
    %   Raw WSL cases: /home/mgnadlinger/Coles_TurbZPGFP/{kwSST,SA}/
    %     20000/wallShearStress and constant/polyMesh, wallViscous patch.
    %     Cfx = abs(tau_x)/q_inf; q_inf = 9572.46608287679 Pa.
    %     Experimental Rex is converted to x with Re/m = 6983542.88149546.
    %   Raw iteration-20000 U and wallShearStress hashes were rechecked against
    %     validation/outputs_nasa_matched_20000/validation_summary.json for both
    %     cases and match. The manuscript PDF is an exact copy of the original
    %     export; the two velocity plots underwent a later typography update.
    % End figure provenance.
    \includegraphics[width=\linewidth]{04_OpenFOAM_HiSA_VnV/Figures/06_Coles_skin_friction.pdf}
    \caption[Coles case: skin-friction compared with NASA and experiment]{%
        Streamwise skin-friction coefficient: full plate
        (top left), leading-edge detail (top right) and downstream detail
        (bottom). Black symbols denote the three archived Coles fence-trip series \cite{Coles1953SupersonicFlatPlate}.}
    \label{fig:coles-skin-friction}
\end{figure}

\paragraph{Assessment.}

With matched physical mesh dimensions, freestream Reynolds number and
molecular and turbulent Prandtl numbers, OpenFOAM/HiSA reproduces the NASA
velocity profiles and downstream wall shear closely.  The result supports the consistency of the turbulence-model implementations for an attached, adiabatic compressible boundary layer.
